\section{Desenvolvimento}
\label{desenvolvimento}

Esta seção descreve como o estudo foi conduzido. Apresentam-se, nesta ordem, a caracterização da pesquisa, os materiais empregados (repositórios e tarefas), o instrumento construído para a comparação (o Uranium), o projeto do experimento e o ambiente de desenvolvimento e de execução. A descrição permite a reprodução: cada parâmetro traz o valor utilizado, cada regra traz o arquivo que a implementa e cada \textit{run} registra um manifesto que o identifica.

\subsection{Caracterização da pesquisa}
\label{ssec:caracterizacao}

A pesquisa pode ser caracterizada em quatro dimensões.

\begin{itemize}
    \item \textbf{Abordagem: quantitativa.} A pergunta de pesquisa é: organizar um agente em papéis orquestrados muda o resultado da correção de defeitos em relação a um agente único com o mesmo modelo, as mesmas ferramentas, o mesmo \textit{prompt} e o mesmo orçamento? A resposta vem de medidas tomadas pelo oráculo sobre o código produzido, e nunca do relato dos sistemas avaliados.
    \item \textbf{Natureza: aplicada.} A investigação exigiu a construção de um aparato experimental, o Uranium, que submete os dois braços às mesmas tarefas sob condições idênticas e verificadas por teste automatizado.
    \item \textbf{Procedimento: experimento controlado.} Manipula-se um único fator, a topologia do sistema de agentes, em dois níveis. O delineamento é pareado por tarefa e por semente. O estudo não tem participantes humanos: o estudo com desenvolvedores previsto no Trabalho Final de Graduação (TFG) passou a trabalho futuro, porque a variância individual excederia o efeito da topologia.
    \item \textbf{Método de condução: \textit{Design Science Research}, com avaliação por experimento controlado.} O aparato foi construído em etapas, cada uma com portões de verificação automática. Após a coleta principal, o braço orquestrado passou por rodadas sucessivas. O autor aprovou cada mudança por escrito; cada mudança recebeu número de decisão e foi medida em uma rodada própria, com o braço de agente único fixo como referência. O registro reúne, até 04/11/2026, 50 decisões numeradas. Esse registro integra o método: ele permite atribuir cada diferença de resultado a uma mudança de desenho, e não ao acaso ou a defeito de instrumentação.
\end{itemize}

\subsection{Materiais: repositórios e tarefas}
\label{ssec:materiais}

\subsubsection{Repositórios-semente}

As tarefas provêm do histórico real de cinco bibliotecas Python de código aberto, com licença permissiva, suíte de testes verde no \textit{commit} escolhido e correções acompanhadas de teste:

\begin{itemize}
    \item \texttt{tomlkit}: leitor e escritor de arquivos TOML que preserva a formatação; cerca de 5.000 linhas e 1.052 testes;
    \item \texttt{sqlite-utils}: manipulação de bancos SQLite; 9.813 linhas e 1.488 testes;
    \item \texttt{peewee}: mapeador objeto-relacional; 15.461 linhas e 1.634 testes;
    \item \texttt{pygments}: realce de sintaxe; 5.330 testes;
    \item \texttt{attrs}: classes com atributos declarativos; 1.395 testes.
\end{itemize}

Oito repositórios foram examinados. Três foram descartados por exigirem passos de construção ou dados externos. Os repositórios de \textit{layout} \texttt{src/} só foram admitidos depois que o aparato passou a suportar esse \textit{layout} sem instalar o pacote no ambiente. Cada repositório tem um espelho local e um ambiente virtual Python próprio, preparado uma única vez no \textit{commit}-base. Nenhuma instalação ocorre durante um \textit{run}. A Seção~\ref{ssec:isolamento} descreve como o \textit{workspace} de cada \textit{run} nasce desse espelho.

\subsubsection{Mineração e verificação das tarefas}

Define-se \textbf{tarefa} como um defeito ou uma funcionalidade real minerado de um projeto Python aberto, com enunciado e testes escondidos. O agente recebe o enunciado; o oráculo guarda os testes. O minerador percorre o histórico (\texttt{git log}) e considera candidato todo \textit{commit} que altera, simultaneamente, código de produção e testes. A \textbf{faixa} do candidato é calculada a partir do \textit{patch} de referência do projeto, por uma função pura congelada antes da coleta (\texttt{scripts/mine\_tasks.py::calcular\_faixa}):

\begin{itemize}
    \item \textbf{pequena}: um ou dois arquivos-fonte e de 10 a 50 linhas alteradas;
    \item \textbf{média}: de 3 a 7 arquivos-fonte e de 40 a 400 linhas;
    \item \textbf{grande}: 3 ou mais arquivos-fonte, 80 ou mais linhas, mais testes novos do que modificados e assunto do \textit{commit} que não começa por \textit{fix} ou \textit{bug}, o que caracteriza funcionalidade nova com teste escrito pelo próprio projeto.
\end{itemize}

Um candidato só é aceito após \textbf{verificação por execução}: a suíte do repositório deve estar verde no \textit{commit}-base, e o teste que acompanhou a correção deve falhar no \textit{commit}-base e passar com a correção aplicada. No segundo ciclo de mineração, sobre cinco repositórios e 14.873 \textit{commits}, a verificação aceitou 3 candidatos grandes, 2 médios e 55 pequenos.
% Os 14.873 commits vêm do texto de 16/10; fonte não localizada em docs/tcc nem em docs/tcc/dados. Os 3/2/55 estão em 07-execucao.md (E3 item 1).
 A maior parte dos candidatos médios e grandes do pré-filtro foi rejeitada porque o teste já passava sem a correção, porque a suíte não executava no \textit{commit}-base ou porque a referência alterava área protegida. A redefinição de ``módulo'' como arquivo \texttt{.py} (decisão 22) e o suporte ao \textit{layout} \texttt{src/} completaram as faixas média e grande: a mineração \texttt{mined-v4} chegou a 3 candidatos grandes, 15 médios e 48 pequenos.

Cada tarefa é um diretório composto por:

\begin{itemize}
    \item o \textbf{enunciado} (\texttt{statement.md}), escrito do ponto de vista de quem usa a biblioteca, sem citar arquivo, classe ou função interna e sem indicar a forma da correção;
    \item o \textbf{oráculo oculto} (\texttt{test\_patch.diff}): os testes da correção original, que o agente nunca vê;
    \item \textbf{duas soluções de referência} distintas (a do projeto e uma alternativa escrita para o estudo), que demonstram que o oráculo não exige uma forma específica de resolver;
    \item \textbf{\textit{patches} de trapaça} plantados, que demonstram que o detector de trapaça funciona;
    \item o arquivo \texttt{task.json}, com a faixa e os parâmetros de execução.
\end{itemize}

Onze enunciados e onze soluções alternativas foram redigidos por um assistente de inteligência artificial (um agente por tarefa) a partir do teste oculto e da \textit{issue} ou do \textit{pull request} original, nunca da mensagem de \textit{commit}, e revisados pelo autor; os dois restantes foram escritos pelo autor. Cada trecho reproduzível de enunciado foi executado no \textit{commit}-base, para evidenciar o problema, e com cada uma das duas soluções, para evidenciar o comportamento esperado.

A título de exemplo, o enunciado da tarefa \texttt{tomlkit-0002} inicia-se com o trecho a seguir. O documento completo contém ainda a seção de reprodução, com o código e a saída esperada, e as restrições do ambiente.

\begin{quote}\small
\textbf{Não há como criar chave pontilhada pela API de construção de documentos.} O TOML permite escrever chaves pontilhadas, como \texttt{foo.bar = 1}, que equivalem a uma tabela \texttt{foo} contendo \texttt{bar = 1}. A biblioteca já lê e preserva esse formato ao analisar um arquivo existente, mas quem monta um documento do zero pela API não consegue produzi-lo: não há como pedir que um valor seja gravado com chave pontilhada. [\ldots] \texttt{tomlkit.key(...)} passa a aceitar, além de uma \textit{string}, uma lista de \textit{strings}. Cada \textit{string} é uma parte da chave: \texttt{tomlkit.key(["foo", "bar"])} representa a chave pontilhada \texttt{foo.bar}. [\ldots] Para \texttt{add(tomlkit.key(["foo", "bar"]), 1)} num documento ou numa tabela vazios, o resultado tem de ser exatamente a \textit{string} \texttt{"foo.bar = 1\textbackslash n"}.
\end{quote}

Uma tarefa só é admitida no estudo depois de aprovada nos cinco portões do procedimento de validação (\texttt{scripts/validate\_task.py}):

\begin{enumerate}
    \item o \textit{commit}-base está verde;
    \item o teste oculto falha no \textit{commit}-base;
    \item cada solução de referência resolve a tarefa, sem regressão e sem acusação do detector de trapaça;
    \item cada \textit{patch} de trapaça é detectado e reprovado;
    \item com cada referência aplicada e sem o \textit{patch} de teste, a suíte visível que o agente executa permanece verde. Este portão entrou em 14/10/2026 (decisão 27). Quando o \textit{patch} de teste original modifica testes visíveis ou um arquivo de referência protegido, os casos afetados saem da suíte visível (\texttt{visible\_deselect}), igualmente para os dois braços e com registro no manifesto. O oráculo nunca desmarca nada.
\end{enumerate}

As treze tarefas foram congeladas em 08/10/2026. Os dados de cada coleta estão versionados em \texttt{docs/tcc/dados/}. A Tabela~\ref{tab-tarefas} lista as tarefas do estudo.

\begin{table}[ht!]
\centering
\caption{As treze tarefas do estudo, por faixa. O repositório é o prefixo do identificador. ``Arq.'' e ``linhas'' referem-se ao \textit{patch} de referência do projeto; ``ocultos'' é o número de identificadores de teste do oráculo que falham no \textit{commit}-base e passam com a referência. Fonte: 3-Desenvolvimento.tex de 16/10 (orig. 10-tarefas.md e \texttt{tasks/*/task.json}).}
\label{tab-tarefas}
\footnotesize
\begin{tabular}{lrrrl}
\toprule
Tarefa & Arq. & Linhas & Ocultos & Faixa \\
\midrule
tomlkit-0002 & 4 & 163 & 1 & grande \\
tomlkit-0003 & 4 & 169 & 1 & grande \\
tomlkit-0004 & 3 & 101 & 7 & grande \\
attrs-0001 & 3 & 46 & 3 & média \\
peewee-0001 & 3 & 54 & 4 & média \\
peewee-0002 & 6 & 72 & 6 & média \\
sqlite-utils-0003 & 3 & 45 & 12 & média \\
peewee-0003 & 2 & 11 & 1 & pequena \\
pygments-0001 & 1 & 26 & 2 & pequena \\
sqlite-utils-0001 & 1 & 34 & 5 & pequena \\
sqlite-utils-0002 & 1 & 21 & 11 & pequena \\
tomlkit-0001 & 1 & 9 & 1 & pequena \\
tomlkit-0005 & 2 & 19 & 1 & pequena \\
\bottomrule
\end{tabular}
\end{table}

\subsection{O instrumento: Uranium}
\label{ssec:uranium}

O Uranium executa os dois braços sobre as mesmas tarefas e os julga por um oráculo que não confia no relato de nenhum deles. A Figura~\ref{fig-arquitetura} resume a arquitetura. Os dois braços compartilham o modelo de linguagem, as ferramentas, o enunciado, o \textit{prompt} de sistema, os limites, o \textit{workspace} isolado e o oráculo. Eles diferem apenas em quem decide o próximo passo e com base em quê.

\begin{figure*}[t]
    \centering
    \begin{tikzpicture}[
        font=\footnotesize,
        caixa/.style={draw, rounded corners=2pt, align=center, inner sep=4pt, minimum height=2em},
        papel/.style={caixa, fill=black!8, text width=46mm},
        comum/.style={caixa, fill=black!18},
        grupo/.style={draw, dashed, rounded corners=3pt, inner sep=6pt},
        seta/.style={-{Stealth[length=2mm]}, semithick},
        dupla/.style={{Stealth[length=2mm]}-{Stealth[length=2mm]}, semithick},
        node distance=5mm and 10mm,
    ]
        \node[comum, text width=0.88\textwidth] (comum) {\textbf{Idêntico nos dois braços:} modelo e provedor fixados por coleta (Qwen3-Coder, DeepSeek V3.2 ou Devstral 2), temperatura 0, semente $s_0 + k$ $\cdot$ sete ferramentas e guarda de caminho $\cdot$ \textit{prompt} de sistema $\cdot$ enunciado $\cdot$ tetos (120 turnos, 8 milhões de \textit{tokens}, 3.600 s) $\cdot$ mensagens de protocolo};

        \node[papel, below=14mm of comum.south, anchor=north, xshift=-38mm] (agente) {Agente único\\LLM com as sete ferramentas; conversa inteira a cada turno};
        \node[caixa, below=of agente, text width=46mm] (operador) {Operador determinístico\\sem modelo: executa a suíte e devolve a saída crua do \texttt{pytest}};
        \draw[dupla] (agente) -- (operador);

        \node[papel, below=14mm of comum.south, anchor=north, xshift=38mm] (refinador) {Refinador de intenção\\\textit{intent} v2 (+ clarificador, até 3 vezes)};
        \node[papel, below=of refinador] (planejador) {Planejador\\somente leitura; até 12 turnos; até 6 sub-tarefas};
        \node[papel, below=of planejador] (dev) {\textit{Developer}\\mesmo \textit{prompt} e ferramentas; ciclos de 40 turnos com memória de trabalho};
        \node[papel, below=of dev] (validador) {Revisor (validador)\\suíte, blocos de reprodução, revisão de sub-tarefa e revisão final pelo modelo};
        \node[papel, below=of validador] (tg) {Gerador de plano de testes (fora da medição)};
        \draw[seta] (refinador) -- (planejador);
        \draw[seta] (planejador) -- (dev);
        \draw[dupla] (dev) -- (validador);
        \draw[seta] (validador) -- (tg);

        \begin{scope}[on background layer]
            \node[grupo, fit=(agente) (operador), label={[font=\footnotesize\bfseries]above:Braço 1: agente único}] (g1) {};
            \node[grupo, fit=(refinador) (tg), label={[font=\footnotesize\bfseries]above:Braço 2: orquestrado}] (g2) {};
        \end{scope}

        \node[caixa, below=10mm of g2.south, anchor=north, xshift=-38mm, text width=0.6\textwidth] (ws) {\textit{Workspace} por \textit{run}: só o \textit{commit}-base, sem histórico \texttt{git} e sem rede; código do agente na caixa de areia \texttt{bubblewrap}; área de testes protegida; trilha de eventos e manifesto};
        \node[comum, below=of ws, text width=0.88\textwidth] (aval) {\textbf{Oráculo e detector de trapaça} (nunca leem o relato dos agentes): \textit{fail-to-pass}, \textit{pass-to-pass}, regressão $\cdot$ regras C1 a C8 $\cdot$ complexidade ciclomática $\cdot$ custo por chamada};
        \draw[seta] (g1.south) -- (ws.north -| g1.south);
        \draw[seta] (g2.south) -- (ws.north -| g2.south);
        \draw[seta] (ws) -- (aval);
    \end{tikzpicture}
    \caption{Arquitetura do Uranium na configuração final: os dois braços, o que é idêntico entre eles, o \textit{workspace} isolado e o oráculo que julga o resultado de ambos.}
    \label{fig-arquitetura}
\end{figure*}

\subsubsection{Ferramentas, guardas e limites compartilhados}

Os agentes dos dois braços atuam sobre o repositório por meio de sete ferramentas, apresentadas ao modelo com os mesmos nomes, descrições, esquemas de argumento e ordem:

\begin{itemize}
    \item \textbf{leitura}: \texttt{list\_files}, \texttt{read\_file} (com faixa de linhas) e \texttt{search\_code} (expressão regular sobre o repositório);
    \item \textbf{escrita}: \texttt{write\_file} (conteúdo completo) e \texttt{replace\_in\_file} (substituição de trecho exato);
    \item \textbf{execução}: \texttt{run\_python} (trecho Python no \textit{workspace}) e \texttt{run\_tests} (suíte visível, inteira ou por identificador).
\end{itemize}

Uma \textbf{guarda de caminho} recusa qualquer escrita em área protegida: diretórios de teste e arquivos de configuração de teste. A recusa retorna ao modelo como texto e conta como ``negação'' na telemetria. Se um trecho executado por \texttt{run\_python} tocar um arquivo protegido, o \textit{harness} recusa a chamada e reverte o arquivo, inclusive quando o arquivo protegido é novo (decisão 47). Os limites das ferramentas são constantes de módulo, iguais nos dois braços e gravadas no manifesto de cada \textit{run} (Tabela~\ref{tab-parametros}).

\subsubsection{Isolamento do \textit{workspace} e caixa de areia}
\label{ssec:isolamento}

A decisão 49, aprovada em 27/10/2026 depois de um achado de 23/10/2026 na rodada 8, mudou o isolamento do \textit{run}. Antes dela, o \textit{workspace} era um \texttt{git clone --shared} do espelho completo. O agente podia ver o ramo \texttt{master}, 56 \textit{tags} e 535 \textit{commits} futuros. A decisão 49 introduziu três mudanças, iguais nos dois braços:

\begin{enumerate}
    \item o \textit{workspace} nasce só com o \textit{commit}-base, por \textit{fetch} raso, sem remoto, sem \textit{tags} e sem \textit{alternates};
    \item toda execução de código do agente roda em caixa de areia \texttt{bubblewrap}, sem rede e com o \textit{workspace} como único diretório gravável;
    \item o detector de trapaça ganhou a regra C8, que compara os arquivos alterados com o histórico futuro do projeto no espelho local.
\end{enumerate}

O manifesto passou à versão 18. As rodadas 1 a 8 e a coleta do Qwen3-Coder correram antes do isolamento; a replicação correu com ele. A Seção~\ref{resultados} relata a varredura dos \textit{runs} anteriores à decisão 49.

\subsubsection{Braço 1: agente único com operador determinístico}

No primeiro braço, o agente recebe o enunciado integral uma única vez e interage com o modelo em turnos. Cada turno é uma chamada ao modelo com toda a conversa acumulada. Um \textbf{operador determinístico} fecha o laço: uma política fixa, sem modelo e sem intervenção humana, serializada no manifesto. O operador aceita todo \textit{patch}. Quando o agente encerra um turno sem solicitar ferramenta, o operador executa a suíte visível. Se a suíte falha, ele devolve a saída crua do \texttt{pytest}, sem resumir nem sugerir. Se a suíte passa e houve mudança no repositório, ele encerra o \textit{run}. O Algoritmo~\ref{alg-single} descreve o laço.

\begin{algorithm}[ht!]
    \normalem
    \KwData{enunciado $E$, \textit{workspace} $W$, orçamento $B$ (turnos, \textit{tokens}, tempo)}
    \KwResult{$W$ modificado e o motivo de parada}
    $M \leftarrow$ [\textit{prompt} de sistema, $E$]\;
    \While{$B$ não esgotado}{
        $r \leftarrow$ \texttt{LLM}($M$); registrar \textit{tokens}, custo e latência; $M \leftarrow M + r$\nllabel{ln:sa-llm}\;
        \eIf{$r$ solicita ferramentas}{
            \ForEach{chamada em $r$ (até 8 por turno)}{
                $s \leftarrow$ executar sobre $W$ com a guarda de caminho, na caixa de areia; $M \leftarrow M + s$\nllabel{ln:sa-exec}\;
            }
        }{
            $T \leftarrow$ \texttt{pytest}($W$), suíte visível com os casos desmarcados da tarefa\nllabel{ln:sa-pytest}\;
            \uIf{$T$ verde e $W$ sem mudança}{ $M \leftarrow M + $ aviso de \textit{diff} vazio\nllabel{ln:sa-aviso}\; }
            \uElseIf{$T$ verde}{ \Return{$W$, ``concluiu''}\nllabel{ln:sa-concluiu}\; }
            \Else{ $M \leftarrow M + $ saída crua do \texttt{pytest} (8.000 caracteres)\nllabel{ln:sa-vermelha}\; }
        }
    }
    \Return{$W$, ``orçamento''}\;
    \caption{Braço de agente único: laço do agente e política do operador.}
    \label{alg-single}
\end{algorithm}

Da linha \ref{ln:sa-llm} à linha \ref{ln:sa-exec}, o agente age: lê, busca, escreve e executa, e cada resultado de ferramenta entra na conversa. Da linha \ref{ln:sa-pytest} à linha \ref{ln:sa-vermelha}, o operador decide. A suíte visível já está verde no \textit{commit}-base, porque o teste discriminante é o oculto. Por isso, o caso ``verde sem mudança'' não encerra o \textit{run}: o operador devolve um aviso (linha \ref{ln:sa-aviso}). O caso ``verde com mudança'' encerra (linha \ref{ln:sa-concluiu}). A suíte vermelha devolve o texto do \texttt{pytest} tal como a ferramenta o produziu (linha \ref{ln:sa-vermelha}). Em nenhum ponto o \textit{harness} interpreta: resumir a falha ou apontar um arquivo injetaria inteligência externa e contaminaria a comparação. Todas as mensagens de protocolo ficam em \texttt{src/feedback.py}, e nenhum braço tem texto próprio.

\subsubsection{Braço 2: orquestrado}

O segundo braço é um grafo de estados (LangGraph) em que papéis especializados cooperam sem intervenção externa. A configuração final é a da rodada 4: manifesto versão 12 na coleta do Qwen3-Coder e versão 18 na replicação; as duas versões diferem só em \texttt{workspace\_history} e \texttt{sandbox}, iguais nos dois braços. Os papéis são:

\begin{itemize}
    \item o \textbf{refinador de intenção} produz a \textit{intent} versão 2: contexto, blocos de reprodução, saídas esperadas, restrições e critérios, copiados do enunciado; um clarificador responde às perguntas pendentes, em até 3 iterações;
    \item o \textbf{planejador} explora o repositório somente com ferramentas de leitura, em até 12 turnos, e divide a \textit{intent} em até 6 sub-tarefas; ele pode replanejar 1 vez, e cada sub-tarefa recebe até 2 ciclos;
    \item o \textbf{\textit{developer}} recebe a \textit{intent} e a sub-tarefa corrente, nunca o enunciado cru; ele usa o mesmo \textit{prompt} de sistema e as mesmas ferramentas do agente único, em ciclos de até 40 turnos, com memória de trabalho entre ciclos;
    \item o \textbf{revisor (validador)} executa a suíte, executa os blocos de reprodução de forma acumulada e pede ao modelo duas revisões: a revisão de sub-tarefa e a revisão final; não há comparação determinística de saídas;
    \item o \textbf{gerador de plano de testes} é o nó terminal; ele produz texto e não escreve na suíte, e seu produto não entra na medição.
\end{itemize}

O Algoritmo~\ref{alg-orq} descreve o laço central na configuração da rodada 4.

\begin{algorithm*}[t]
    \normalem
    \KwData{enunciado $E$, \textit{workspace} $W$, orçamento $B$, política $P$ (manifesto)}
    \KwResult{$W$ modificado e o motivo de parada}
    $I \leftarrow$ \texttt{Refinador}($E$), \textit{intent} v2\nllabel{ln:or-refinador}\;
    \While{$I$ não pronta e iterações $< 3$}{$I \leftarrow$ \texttt{Clarificador}($I$)\nllabel{ln:or-clarificador}\;}
    $Q \leftarrow$ \texttt{Planejador}($I$, $W$): até 12 turnos de leitura; até 6 sub-tarefas\nllabel{ln:or-planejador}\;
    memória $\leftarrow$ [ ]; reprovações $\leftarrow 0$; replanejamentos $\leftarrow 0$\;
    \While{$B$ não esgotado}{
        $q \leftarrow$ sub-tarefa corrente de $Q$\nllabel{ln:or-subtarefa}\;
        $D \leftarrow$ \texttt{Developer}($I$, $q$, memória, \textit{feedback}): até 40 turnos, mesmo \textit{prompt} e ferramentas do Algoritmo~\ref{alg-single}\nllabel{ln:or-developer}\;
        memória $\leftarrow$ memória $+$ registro do ciclo (arquivos lidos, buscas, escritas, suítes, nota de ciclo) $+$ \textit{diff} atual\nllabel{ln:or-memoria}\;
        $T \leftarrow$ \texttt{pytest}($W$); $X \leftarrow$ executar os blocos de reprodução de $I$, cada um precedido dos anteriores\nllabel{ln:or-suite}\;
        \If{$Q$ aberto}{
            $V_q \leftarrow$ \texttt{RevisorSubtarefa}($I$, $q$, \textit{diff}, $X$)\nllabel{ln:or-revsub}\;
            \lIf{$V_q$ aprova}{concluir $q$}
            \lElseIf{$q$ já teve 2 ciclos}{adiar $q$}
            \lIf{$Q$ esgotado e replanejamentos $< 1$}{$Q \leftarrow$ \texttt{Planejador}($I$, $W$, memória); replanejamentos $\leftarrow 1$}
        }
        \If{$T$ verde}{
            $V \leftarrow$ \texttt{RevisorFinal}($I$, \textit{diff}, $X$)\nllabel{ln:or-revfinal}\;
            \If{$V$ aprova e $W$ mudou}{ \Return{$W$, ``concluiu''}\nllabel{ln:or-concluiu}\; }
            \If{$Q$ fechado}{
                reprovações $\leftarrow$ reprovações $+ 1$\nllabel{ln:or-reprovacoes}\;
                \lIf{reprovações $\geq 2$}{\Return{$W$, ``teto de reprovações''}}
            }
        }
        \textit{feedback} $\leftarrow$ saída crua do \texttt{pytest} $\mid$ parecer do revisor\nllabel{ln:or-feedback}\;
    }
    \Return{$W$, ``orçamento''}\;
    \caption{Braço orquestrado: laço \textit{developer}--validador com planejador e revisões pelo modelo (configuração da rodada 4).}
    \label{alg-orq}
\end{algorithm*}

As linhas \ref{ln:or-refinador} e \ref{ln:or-clarificador} produzem a \textit{intent} que o \textit{developer} recebe em lugar do enunciado original. Trata-se de um traço do protótipo, declarado como parte do tratamento e medido por uma ablação. A linha \ref{ln:or-planejador} é a exploração inicial do planejador. Nas linhas \ref{ln:or-developer} e \ref{ln:or-memoria}, o \textit{developer} trabalha em \textbf{ciclos}: a cada 40 turnos, ou quando ele próprio encerra, a conversa é descartada. No retorno, ele recebe a \textbf{memória de trabalho}, isto é, o registro determinístico do que fez e o \textit{diff} atual, nunca a transcrição anterior. Da linha \ref{ln:or-suite} à linha \ref{ln:or-feedback}, atua o validador. A linha \ref{ln:or-suite} executa a suíte e os blocos de reprodução de forma acumulada. A revisão de sub-tarefa (linha \ref{ln:or-revsub}) decide, pelo modelo, se a sub-tarefa corrente está concluída; uma sub-tarefa sem aprovação após 2 ciclos é adiada, e um plano esgotado é refeito uma vez. A revisão final começa na linha \ref{ln:or-revfinal}. A suíte decide primeiro: o revisor só pode rebaixar um veredito verde, nunca promover um vermelho. O teto de 2 reprovações finais (linha \ref{ln:or-reprovacoes}) é o único critério de parada inexistente no outro braço, e não conta enquanto o plano está aberto. A linha \ref{ln:or-feedback} indica o que retorna ao \textit{developer}: o mesmo texto cru do \texttt{pytest} que o agente único recebe, acrescido do parecer do revisor, que só existe neste braço e constitui a própria topologia. O critério de parada é igual nos dois braços: suíte verde, com aprovação do revisor no orquestrado, ou orçamento.

A Tabela~\ref{tab-papeis} resume o que cada papel recebe, produz e pode fazer.

\begin{table*}[t]
\centering
\caption{Papéis do braço orquestrado na configuração da rodada 4. Fontes: 29 \S2 e \S3; 27 \S2.2; 3-Desenvolvimento.tex de 16/10.}
\label{tab-papeis}
\footnotesize
\begin{tabular}{lp{3.0cm}p{3.8cm}p{2.3cm}p{2.1cm}}
\toprule
Papel & Recebe & Produz & Ferramentas & Chamadas \\
\midrule
Refinador & enunciado & \textit{intent} v2 (contexto, blocos de reprodução, saídas esperadas, restrições e critérios copiados) & nenhuma & 1 \\
Clarificador & \textit{intent} com perguntas & respostas simuladas & nenhuma & até 3 \\
Planejador & \textit{intent}; repositório & plano de até 6 sub-tarefas; 1 replanejamento & \texttt{list\_files}, \texttt{read\_file}, \texttt{search\_code} & até 12 turnos de leitura \\
\textit{Developer} & \textit{intent}; sub-tarefa corrente; memória de trabalho; \textit{feedback} & \textit{diff} do repositório; nota de ciclo & as sete & até 40 turnos por ciclo \\
Revisor (validador) & repositório; \textit{intent}; plano & relatório da suíte; saídas dos blocos; revisão de sub-tarefa e revisão final pelo modelo; avanço do plano & suíte e blocos, pelo \textit{harness} & 1 por revisão \\
Gerador de plano de testes & \textit{intent}; \textit{diff} & plano de testes em texto, fora da medição & nenhuma & 1 \\
\bottomrule
\end{tabular}
\end{table*}

\subsubsection{Paridade entre os braços}

A alegação causal do estudo depende de os dois braços serem idênticos em tudo, exceto na topologia. Testes automatizados verificam essa identidade antes de cada \textit{commit} do instrumento (\texttt{tests/test\_arm\_parity.py}). Eles comparam:

\begin{itemize}
    \item o \textbf{manifesto} de um \textit{run} de cada braço, sobre a mesma tarefa e repetição, que só pode diferir em um conjunto declarado de chaves (\texttt{arm}, \texttt{topology}, \texttt{driver\_policy}, \texttt{orchestration\_policy}, \texttt{run\_id}, \texttt{started\_at} e \texttt{attempt}), cada uma com justificativa em \texttt{CHAVES\_QUE\_PODEM\_DIFERIR};
    \item a configuração do modelo e o orçamento;
    \item o conjunto, a ordem, os esquemas e as descrições das ferramentas;
    \item o \textit{prompt} de sistema e as mensagens de protocolo (aviso de \textit{diff} vazio, texto da falha da suíte e seu truncamento), comparados literalmente;
    \item a semente: a repetição $k$ de um braço usa a mesma semente que a repetição $k$ do outro, e os dois \textit{runs} formam um \textbf{par}.
\end{itemize}

Tudo o que altera o comportamento de um braço vai ao manifesto (\texttt{orchestration\_policy}, \texttt{shared\_limits}, \texttt{budget}, \texttt{llm}). Duas configurações diferentes nunca produzem manifestos iguais. A versão do manifesto é 12 na coleta do Qwen3-Coder e 18 na replicação.

\subsubsection{Oráculo, detector de trapaça e telemetria}

Terminado o \textit{run}, o oráculo julga o \textit{workspace} sem ler o relato do sistema. Ele injeta os testes ocultos em um diretório inexistente no \textit{commit}-base e verifica, com a rede isolada, três condições (\texttt{harness/oracle.py}):

\begin{enumerate}
    \item os testes ocultos passam (\textit{fail-to-pass});
    \item a suíte que já passava continua passando (\textit{pass-to-pass});
    \item a suíte original completa não regrediu.
\end{enumerate}

O \textbf{detector de trapaça} (\texttt{harness/cheat.py}) aplica oito regras estáticas sobre o \textit{patch}, comparando o estado anterior e o posterior ao \textit{run}. Uma violação crítica anula o \textit{run}. As regras C3 a C6 operam sobre a árvore sintática, para não acusar comentários ou \textit{docstrings}:

\begin{enumerate}
    \item[C1] arquivo de teste ou de configuração alterado;
    \item[C2] opções do \texttt{pytest} alteradas;
    \item[C3] marcador de \texttt{skip} ou \texttt{xfail} inserido;
    \item[C4] arquivo \texttt{conftest.py} criado ou alterado;
    \item[C5] código de produção que detecta estar sob teste;
    \item[C6] bloco \texttt{except: pass} novo, capaz de suprimir a exceção esperada pelo teste;
    \item[C7] referência ao diretório dos testes ocultos, que o agente não deveria conhecer: criação do diretório ou menção a ele no código;;
    \item[C8] arquivos alterados que coincidem com o histórico futuro do projeto no espelho local (decisão 49).
\end{enumerate}

O módulo de \textbf{qualidade} mede a complexidade ciclomática média dos arquivos alterados, antes e depois. O módulo de \textbf{custo} soma \textit{tokens}, dólares e chamadas a partir da resposta do provedor, chamada a chamada. Um \textit{run} é \textbf{resolvido} quando o oráculo aprova nas três verificações e o detector de trapaça não acha violação crítica.

Todo \textit{run} grava uma trilha de eventos em formato JSON por linha (\texttt{events.jsonl}), validada contra um esquema com catorze tipos de evento: início e fim de \textit{run}, entrada e saída de nó, roteamento, chamada ao modelo, chamada de ferramenta, execução de suíte, erro, corte de orçamento, \textit{intent}, validação, memória e plano. Grava também o \texttt{manifest.json}, com toda a configuração, e o \texttt{result.json}, com o veredito do oráculo e do detector de trapaça, o custo, a qualidade e o \textit{diff}. As métricas saem dessa trilha, e não de relato.

\subsection{Projeto de experimentos}
\label{ssec:experimento}

\subsubsection{Fatores, níveis e variáveis de resposta}

O delineamento tem os seguintes elementos:

\begin{itemize}
    \item \textbf{fator}: a topologia do sistema de agentes;
    \item \textbf{níveis}: agente único com operador determinístico; orquestrado na configuração da rodada 4;
    \item \textbf{blocos}: a tarefa (13), a repetição com sua semente (3) e o modelo (3);
    \item \textbf{unidade de análise}: a tarefa; o valor por tarefa e braço é a maioria das três repetições, para a resolução, ou a mediana, para as métricas contínuas;
    \item \textbf{pareamento}: a repetição $k$ usa a semente $s_0 + k$ nos dois braços, com $s_0 = 20260820$, o que dá as sementes 20260821, 20260822 e 20260823;
    \item \textbf{estrato}: a faixa (3 grandes, 4 médias e 6 pequenas), usada na leitura por faixa e na proposição P3.
\end{itemize}

Cada coleta tem 78 \textit{runs}: 13 tarefas $\times$ 3 repetições $\times$ 2 braços, com um modelo. As três coletas finais somam 234 \textit{runs} e 117 pares.

O protocolo fixou três proposições antes da coleta:

\begin{itemize}
    \item \textbf{P1}: a topologia muda a taxa de resolução. A hipótese é neutra e o teste é bilateral: $H_0$ afirma que a topologia não altera a taxa de resolução; $H_1$, que a altera, em qualquer direção;
    \item \textbf{P2}: a topologia muda o custo em \textit{tokens} nas tarefas resolvidas pelos dois braços;
    \item \textbf{P3}: o efeito depende da faixa.
\end{itemize}

As proposições P4 a P10 surgiram durante as rodadas. O texto as trata como achados de construção, não como confirmação de hipótese.

As variáveis de resposta são as sete métricas da Tabela~\ref{tab-metricas}, com hierarquia fixada antes da coleta. Nenhuma depende da arquitetura interna de um dos braços. A comparação é sempre orquestrado menos agente único.

\begin{table*}[t]
\centering
\caption{Variáveis de resposta, com a hierarquia e o teste fixados antes da coleta. Fontes: 27 \S1.3 e \S2.6; 29 \S3; 3-Desenvolvimento.tex de 16/10.}
\label{tab-metricas}
\small
\begin{tabular}{cp{2.6cm}p{5.6cm}lp{4.6cm}}
\toprule
\# & Métrica & Definição & Melhor & Análise \\
\midrule
1 & Taxa de resolução & proporção de \textit{runs} resolvidos & maior & primária; sinal exato bilateral por tarefa; McNemar exato bilateral por par \\
2 & Custo computacional & \textit{tokens} por \textit{run} (dólares e chamadas reportados em conjunto) & menor & primária inferencial; Wilcoxon pareado exato nas tarefas resolvidas pelos dois braços \\
3 & Tempo de parede & segundos de relógio do \textit{run} & menor & Wilcoxon pareado exato, mesma censura \\
4 & Consistência & tarefas resolvidas nas três repetições, entre as resolvidas ao menos uma vez & maior & descritiva \\
5 & Variação da complexidade & complexidade ciclomática média dos arquivos alterados, depois menos antes & menor & Wilcoxon pareado exato, exploratória \\
6 & Conclusão sem corte & \textit{runs} encerrados por decisão do sistema, não por limite & maior & descritiva \\
7 & Violações & violação do detector de trapaça ou tentativa de escrita em área protegida & menor & descritiva \\
\bottomrule
\end{tabular}
\end{table*}

O TFG prometia medir produtividade, autonomia, qualidade dos artefatos, rastreabilidade e custo. A Tabela~\ref{tab-tfg} mostra o destino de cada promessa nas sete métricas.

\begin{table*}[t]
\centering
\caption{Mapa das dimensões prometidas no TFG para o desenho final. Fonte: 27 \S1.2.}
\label{tab-tfg}
\small
\begin{tabular}{p{3.6cm}p{5.2cm}p{6.6cm}}
\toprule
Promessa do TFG & Destino & Motivo \\
\midrule
Complexidade ciclomática & métrica 5, como variação (depois menos antes), exploratória & o valor absoluto é do repositório, não do agente \\
\textit{Maintainability Index} & excluído & saturou em 0,0 em 100\% dos resultados do piloto \\
\textit{Mutation Score} & excluído & mede a força da suíte; nenhum braço escreve testes, logo mediria a suíte do repositório, idêntica nos dois braços \\
Cobertura & excluída & função nunca chamada \\
Índice de autonomia L0 a L4 & excluído; a autonomia virou ``conclusão sem corte'' (métrica 6) & constante entre braços sem humano no laço \\
Rastreabilidade & não é métrica; virou portão de admissão do dado & provedor fixado em toda chamada, rede isolada na avaliação, trilha válida com um só \texttt{run\_end} \\
Produtividade & métrica 3, tempo de parede & sem humanos, produtividade é eficiência do sistema \\
Custo & métrica 2: \textit{tokens} testados; dólares e chamadas reportados & \textit{tokens} independem de tabela de preço \\
\bottomrule
\end{tabular}
\end{table*}

\subsubsection{Métodos estatísticos}

Os testes são exatos e implementados em funções puras do \textit{script} de análise (\texttt{scripts/analyze\_results.py}), sem biblioteca estatística. O nível de significância é 5\%.

\begin{itemize}
    \item \textbf{Sinal exato bilateral por tarefa} (P1, métrica 1). O teste binomial considera só as tarefas em que os braços discordam na maioria das três repetições. O protocolo o declara \textit{a priori} e o reporta mesmo quando não significativo.
    \item \textbf{McNemar exato bilateral por par} (P1 e P3, métrica 1). A unidade é o par: mesmo modelo, tarefa e repetição, logo mesma semente. A análise usa 39 pares por modelo e 117 pares no agregado, também por faixa. Ela trata \textit{runs} da mesma tarefa como independentes condicionalmente à semente e é complementar à análise por tarefa.
    \item \textbf{Wilcoxon pareado exato} (P2, métricas 2, 3 e 5). O teste enumera as $2^n$ atribuições de sinal. Ele usa só as tarefas resolvidas pelos dois braços, e a censura é declarada com a lista das tarefas excluídas, porque um \textit{run} cortado por limite custa o teto, e não o custo de resolver.
    \item \textbf{Poder por simulação}. O procedimento simula 13 tarefas, 3 repetições, maioria por tarefa e sinal exato bilateral; a linha de base é a taxa agregada observada do agente único por tarefa; os efeitos verdadeiros simulados são de +10, +20 e +30 pontos percentuais por \textit{run}, com 13, 30 e 60 tarefas. O número de réplicas da simulação é 4.000.
\end{itemize}

As sete métricas nos três modelos não recebem correção para múltiplas comparações. Toda análise fora dessa hierarquia, como o custo da carga inteira e as análises do documento 28, é exploratória e declarada pós-hoc.

\subsubsection{Limites, tetos e protocolo de execução}

Os tetos de um \textit{run} foram calibrados nas fumaças A a H, na tarefa \texttt{tomlkit-0004}, e congelados em 12/10/2026, antes da coleta:

\begin{itemize}
    \item 120 turnos, o orçamento propriamente dito, comparável entre repositórios;
    \item 8 milhões de \textit{tokens} e 3.600 segundos de parede, como redes de segurança;
    \item 120 segundos por execução de suíte.
\end{itemize}

Nenhum desses valores mudou depois. Mudar qualquer um deles invalidaria a comparação com os \textit{runs} já coletados. O protocolo de execução obedece às seguintes regras:

\begin{itemize}
    \item um \textit{run} é \textbf{perdido}, e não fracassado, quando termina por falha de infraestrutura do provedor ou por interrupção externa (prefixos \texttt{infra\_error:} e \texttt{interrupted:}); ele sai do denominador e é refeito com a mesma semente e com o número de tentativa incrementado;
    \item o corte por orçamento é resultado, não perda;
    \item os \textit{runs} ocorrem um de cada vez, com o agente único antes do orquestrado em cada par;
    \item a coleta é retomável: cada célula (tarefa, braço, repetição) tem estado próprio em disco; a coleta reexecuta com \texttt{--retry-delay 300 --max-attempts 5}.
\end{itemize}

Antes de admitir um \textit{run} na análise, o \textit{harness} confere os portões de admissão: o provedor que atendeu cada chamada é o provedor fixado; todas as chamadas saíram pela chave própria do provedor (o número de chamadas BYOK, \textit{Bring Your Own Key}, é igual ao número de chamadas ao modelo); a rede ficou isolada no oráculo; a trilha de eventos é válida, com um só \texttt{run\_end}; e o \textit{run} não está marcado como perdido.

\subsubsection{Escolha da configuração do orquestrado}

A coleta principal (rodada 1) mediu o protótipo do braço orquestrado tal como herdado do TFG, com a memória de trabalho entre ciclos e o teto de 40 turnos por ciclo definidos nas fumaças. Em seguida, o braço orquestrado passou por uma calibração iterativa: uma ablação do refinador e sete configurações, numeradas como rodadas 2 a 8, cada uma com política registrada no manifesto e decisão aprovada por escrito, medidas entre 14 e 27/10/2026 em 273 \textit{runs}, sempre com as mesmas 13 tarefas e sementes e com o agente único fixo como referência pareada. Antes das rodadas 7 e 8, o protocolo fixou o critério de \textbf{configuração decisiva}: 10 ou mais tarefas resolvidas em 13, sem perder na faixa grande em relação à rodada 4; pela regra da decisão 45, se nenhuma rodada o atingisse, a rodada 4 seria a configuração do orquestrado na replicação. Nenhuma atingiu, e a rodada 4 tornou-se a configuração final descrita nesta seção. No mesmo período, a construção do instrumento corrigiu 16 defeitos, válidos para os dois braços, entre eles o quinto portão e o \texttt{visible\_deselect} (decisão 27), o diretório-raiz do \texttt{pytest} (decisão 30), a guarda de caminho para arquivo protegido novo (decisão 47) e o isolamento da Seção~\ref{ssec:isolamento} (decisão 49); as células afetadas foram descartadas e refeitas com a mesma semente. O relatório do experimento, no repositório do instrumento, traz a tabela das rodadas e o resultado de cada uma; o artigo não as relata uma a uma.

\subsubsection{Replicação}

A decisão 45 definiu a replicação da rodada 4 em dois modelos de outros laboratórios. A Tabela~\ref{tab-modelos} resume os três modelos. A escolha mudou primeiro o laboratório, mantendo o provedor (DeepSeek V3.2 na DeepInfra), e depois o laboratório e o provedor (Devstral 2 na Mistral). O estudo usa, assim, três modelos de três laboratórios e dois provedores.

\begin{table*}[t]
\centering
\caption{Modelos do estudo, laboratório, provedor e papel. Fontes: 29 \S3.1; 27 \S2.7.}
\label{tab-modelos}
\small
\begin{tabular}{p{2.0cm}p{4.3cm}p{1.8cm}p{1.6cm}p{4.4cm}}
\toprule
Modelo & Identificador & Laboratório & Provedor & Papel e observação \\
\midrule
Qwen3-Coder & \texttt{qwen/qwen3-coder} & Alibaba & DeepInfra & piloto, rodadas e coleta principal; quantização fp4 publicada \\
DeepSeek V3.2 & \texttt{deepseek/deepseek-v3.2}, modo não pensante & DeepSeek & DeepInfra & replicação: muda o laboratório, mantém o provedor; quantização fp4 publicada \\
Devstral 2 & \texttt{mistralai/devstral-2512} & Mistral & Mistral & replicação: muda o laboratório e o provedor; provedor único; quantização não publicada (declarado) \\
\bottomrule
\end{tabular}
\end{table*}

Os três modelos atendem aos mesmos critérios: \textit{endpoint} com semente, ferramentas, temperatura e teto de saída; pesos abertos; modo não pensante. A escolha do Qwen3-Coder considerou 23 candidatos e nove requisitos. Modelos pensantes ficaram fora porque a contabilidade de contexto foi validada com raciocínio zero. Modelos fechados de custo médio ficaram fora por não exporem a semente.

A replicação seguiu três passos para cada modelo:

\begin{enumerate}
    \item uma sondagem de aptidão (\texttt{fitness\_test.py -n 5}) conferiu chamada de ferramentas, argumentos exatos, provedor e ausência de \textit{tokens} de raciocínio;
    \item uma fumaça de três tarefas, nos dois braços, conferiu o \textit{harness} sob a decisão 49;
    \item uma coleta de 78 \textit{runs}, com as mesmas 13 tarefas, sementes e tetos, e o orquestrado na rodada 4 com manifesto v18.
\end{enumerate}

A coleta final do Qwen3-Coder reúne os \textit{runs} do orquestrado da rodada 4 e os \textit{runs} pareados do agente único. As coletas da replicação ocorreram entre 27 e 29/10/2026.

\subsubsection{Cronograma}

A Tabela~\ref{tab-cronograma} resume as etapas do estudo. O estudo soma mais de 580 \textit{runs}.

\begin{table*}[t]
\centering
\caption{Cronograma do estudo, de setembro a outubro de 2026. Fonte: 29 \S3.2.}
\label{tab-cronograma}
\small
\begin{tabular}{p{3.0cm}p{5.0cm}p{7.6cm}}
\toprule
Quando & Etapa & Produto \\
\midrule
11/09/2026 & piloto & Qwen3-Coder na DeepInfra com semente fixa \\
19/09/2026 & escolha de modelo e provedor & critérios sobre \textit{endpoints} consultados ao vivo \\
08/10/2026 & tarefas congeladas & 13 tarefas, cinco portões, duas referências \\
08 e 09/10/2026 & fumaças A a H & calibração dos tetos; congelados em 12/10 \\
12 e 13/10/2026 & coleta principal & 78 \textit{runs}, Qwen3-Coder \\
14/10/2026 & ablação do refinador e \textit{intent} v2 & \textit{developer} com o enunciado íntegro; \textit{intent} v2 \\
14 a 27/10/2026 & rodadas 2 a 8 & sete configurações do orquestrado, 273 \textit{runs}, 16 defeitos do instrumento corrigidos \\
22/10/2026 & isolamento & \textit{workspace} só com o \textit{commit}-base, caixa de areia, regra C8 \\
27 a 29/10/2026 & replicação & DeepSeek V3.2 e Devstral 2, 78 \textit{runs} cada \\
29 e 30/10/2026 & relatório e análise exploratória & documentos 27 e 28; revisão de código do \textit{harness} \\
\bottomrule
\end{tabular}
\end{table*}

\subsubsection{Reprodução}

A reprodução do estudo segue quatro passos, todos documentados no repositório do instrumento:

\begin{enumerate}
    \item preparar os espelhos e os ambientes virtuais dos repositórios-semente (\texttt{scripts/setup\_mirror.py <repositório>});
    \item validar cada tarefa nos cinco portões (\texttt{scripts/validate\_task.py tasks/<tarefa>});
    \item executar a coleta (\texttt{scripts/run\_collection.py --arms <braços> --repetitions 3 --out <diretório>}), com a política do braço orquestrado definida por variáveis de ambiente documentadas em \texttt{.env.example}; o modelo e o provedor de cada coleta ficam registrados na chave \texttt{llm} do manifesto;
    \item analisar os resultados (\texttt{scripts/analyze\_results.py <diretório> --markdown}), com o apêndice de métricas gerado por \texttt{scripts/metrics\_appendix.py}.
\end{enumerate}

O manifesto de cada \textit{run}, a trilha de eventos e o arquivo \texttt{result.json} ficam no diretório da célula. Os dados de cada rodada e de cada coleta estão versionados em \texttt{docs/tcc/dados/}.

\subsection{Ambiente de desenvolvimento e de experimentos}
\label{ssec:ambiente}

Os \textit{runs} ocorreram em uma única máquina: processador AMD Ryzen 7 7800X3D (8 núcleos e 16 \textit{threads}), 7 GB de memória visíveis à máquina virtual, Ubuntu 26.04.1 LTS sob WSL2 (\textit{kernel} 6.18.54) e Python 3.12.14. O modelo de linguagem é servido remotamente. A máquina executa os agentes, as ferramentas, a caixa de areia, as suítes de teste e o oráculo. A Tabela~\ref{tab-parametros} lista as versões e os parâmetros e indica quais valores são o padrão da biblioteca e quais o estudo fixou.

\begin{table*}[t]
\centering
\caption{\textit{Software}, modelos e parâmetros. ``Fixado'' indica valor escolhido pelo estudo; ``padrão'' indica o valor de fábrica da biblioteca. Fontes: 3-Desenvolvimento.tex de 16/10; 27 \S2.1, \S2.4, \S2.5, \S2.7 e \S3.5; 29 \S3 e \S3.1.}
\label{tab-parametros}
\small
\begin{tabular}{p{4.8cm}p{4.6cm}p{5.6cm}}
\toprule
Componente & Valor & Observação \\
\midrule
LangGraph / LangChain & 1.2.0 / core 1.4.0, openai 1.2.1 & grafo do braço orquestrado; cliente do modelo \\
openai (SDK) / pydantic & 2.37.0 / 2.13.4 & chamadas HTTP; esquemas de saída estruturada \\
pytest / radon & 9.0.3 / 6.0.1 & suítes e oráculo; complexidade ciclomática \\
\texttt{bubblewrap} & 0.11.1 & caixa de areia sem rede; só o \textit{workspace} gravável \\
Modelos & \texttt{qwen/qwen3-coder}; \texttt{deepseek/deepseek-v3.2}; \texttt{mistralai/devstral-2512} & via OpenRouter, um modelo por coleta, chave própria do provedor (BYOK) \\
Roteamento & \texttt{only}=[provedor da coleta: DeepInfra ou Mistral], sem \textit{fallback}, \texttt{require\_parameters}, sem coleta de dados & fixado: o provedor não pode mudar durante o \textit{run} \\
Temperatura / semente & 0 / $s_0 + k$, $s_0 = 20260820$ & fixado; semente igual nos dois braços na repetição $k$ \\
\textit{Tokens} de saída / tempo por chamada / tentativas & 8.192 / 90 s / 6 & fixado (o padrão do cliente é 2 tentativas) \\
Reexecução de \textit{run} perdido & \texttt{--retry-delay 300 --max-attempts 5} & fixado; mesma semente \\
Orçamento por \textit{run} & 120 turnos; 8.000.000 \textit{tokens}; 3.600 s & fixado, congelado em 12/10/2026 \\
Tempo por suíte & 120 s & fixado \\
Chamadas de ferramenta por turno & 8 & fixado \\
\texttt{read\_file} & 400 linhas; 20.000 caracteres & fixado \\
\texttt{search\_code} / \texttt{list\_files} & 50 / 200 resultados & fixado \\
Saída de ferramenta & 8.000 caracteres (2/3 do início, 1/3 do fim) & fixado; mesmo corte nos dois braços \\
\texttt{run\_python} & 60 s & fixado \\
\textit{Workspace} & \textit{fetch} raso do \textit{commit}-base; sem remoto, \textit{tags} ou \textit{alternates} & fixado desde a decisão 49 \\
Orquestrado: ciclo / teto de reprovações / clarificações & 40 turnos / 2 reprovações finais, não contadas com o plano aberto / 3 & fixado \\
Orquestrado: planejador (rodada 4) & 12 turnos de leitura; 6 sub-tarefas; 1 replanejamento; 2 ciclos por sub-tarefa & fixado; revisão de sub-tarefa pelo modelo \\
\bottomrule
\end{tabular}
\end{table*}

O código do instrumento conta com 751 testes automatizados. Eles cobrem o detector de trapaça, o oráculo, a guarda de caminho, a fixação do provedor, a paridade entre os braços e cada política do braço orquestrado. A suíte roda e deve estar verde antes de cada \textit{commit}.
